\documentclass[11pt,a4paper]{article}
\usepackage[margin=1in]{geometry}
\usepackage{amsmath,amssymb,amsthm,booktabs}
\usepackage[hyphens]{url}
\usepackage[hidelinks]{hyperref}
\newtheorem{theorem}{Theorem}
\newtheorem{corollary}[theorem]{Corollary}
\newcommand{\Z}{\mathbb Z}
\newcommand{\A}{\mathcal A}
\newcommand{\D}{\mathcal D}
\newcommand{\Sset}{\mathcal S}
\title{An explicit weighted digit construction\\for small sumsets and large difference sets}
\author{Benjamin Guo\\Hraness (\url{hraness.com})\thanks{Developed with AI research agents. The proof is accompanied by a separate verifier using exact integer arithmetic.}}
\date{28 September 2026}
\hypersetup{pdftitle={An explicit weighted digit construction for small sumsets and large difference sets},pdfauthor={Benjamin Guo}}

\begin{document}
\maketitle
\begin{abstract}
The small-sumset, large-difference-set problem asks how large $|E-F|$ can be
when $|E+F|\le K|E|$, for fixed $K>1$ and arbitrarily large finite integer
sets $E,F$. We give a weighted digit construction for its exponent $\theta$.
Equal counts of opposite differences give a lower bound for the difference
set, while a product-weight estimate bounds the sumset from above. The
resulting formula applies to any finite digit alphabet. For the alphabet
$\{0,3,4,6,7,\ldots,16\}$ and geometric weights $(16/19)^a$, it proves
$\theta>1.1855$. The numerical conclusion reduces to one explicit integer
inequality. This exceeds Zheng's published lower bound of $1.173077$.
The proof uses the finite-set transfer
inequality of Gyarmati, Hennecart and Ruzsa and does not require optimality
of the chosen parameters.
\end{abstract}
\medskip\noindent\textbf{Keywords.} sumset; difference set; weighted digits; entropy; exact certificate.
\smallskip\noindent\textbf{MSC 2020.} 11B75; 05A16.

\section{The exponent and the finite-set reduction}
Let $\theta$ be the supremum of the real numbers $t$ with the following
property: for every fixed $K>1$, there is a constant $c(K)>0$ and arbitrarily
large finite sets $E,F\subset\Z$ satisfying
\begin{equation}\label{eq:problem}
 |E+F|\le K|E|,
 \qquad
 |E-F|\ge c(K)|E+F|^t.
\end{equation}
Here $E\pm F=\{e\pm f:e\in E,\ f\in F\}$, and all logarithms are natural.
The constant may depend on the fixed exponent $t$ as well as on $K$.

The lemma in Section~2 of Gyarmati, Hennecart and Ruzsa~\cite{ghr} gives
a finite-set transfer inequality: for every finite $U\subset\Z_{\ge0}$
with $0\in U$, $\max U>0$ and $|U-U|<2\max U+1$,
\begin{equation}\label{eq:transfer}
 \theta\ge
 1+\frac{\log\bigl(|U-U|/|U+U|\bigr)}{\log(2\max U+1)}.
\end{equation}
For each such $U$, its exponent is attained for every fixed $K>1$ and
arbitrarily large sets. The strict difference-count hypothesis matters;
we will ensure it by dilating each encoded set. A family whose
difference-to-sum ratio grows exponentially relative to its diameter
then gives a lower bound for the supremum $\theta$.

Zheng uses integer vectors whose coordinates lie in a full interval and
whose coordinate sum is bounded. His large-deviation estimates give the
lower bound $1.173077$~\cite{zheng}. We allow a smaller coordinate alphabet
and general nonnegative digit costs. Matching the counts of a difference
and its negative makes the costs of its two representatives exactly equal.
This symmetry gives a direct counting argument without optimizing a
large-deviation rate function. The constructions below improve the bound
in Zheng's stated theorem.

A separate problem concerns the one-set quantity
\[
 C(A)=\frac{\log(|A+A|/|A|)}{\log(|A-A|/|A|)},
 \qquad A\subset\Z,\quad |A|>1,
\]
whose supremum is $2$ by Lin and Li~\cite{linli}. Numerical comparisons
must distinguish $C(A)$ from $\theta$. The ``Sets'' comparison in
Table~1 of Beck, Ogren and Kobren~\cite{hill} reports a value of $1.21$
from Hyra, while its Appendix~A, equation~(7), and Appendix~D use the
objective considered here. Lin and Li~\cite[Section~3]{linli} identify
that $1.21$ search value with $C(A)$. The comparison therefore does not
establish a $1.21$ lower bound for $\theta$.

\section{A weighted digit construction}
Fix an integer $B\ge1$ and a digit alphabet
$\A\subseteq\{0,\ldots,B\}$ containing at least two elements. Attach a
positive weight $w_a$ to each $a\in\A$. Multiplying every weight by a
common factor changes neither of the ratios below, so assume
$\max_{a\in\A}w_a=1$. Put
\begin{align}
 \Sset&=\A+\A, &
 S_s&=\max_{a+b=s} w_aw_b, &
 P&=\sum_{s\in\Sset}S_s,\label{eq:sums}\\
 \D&=\A-\A, &
 D_d&=\max_{a-b=d} w_aw_b, &
 Q&=\sum_{d\in\D}D_d.\label{eq:differences}
\end{align}
All pairs in these maxima belong to $\A^2$; only represented sums and
differences are included.

\begin{theorem}\label{thm:weighted}
If $Q>P$, then
\begin{equation}\label{eq:bound}
 \theta\ge 1+\frac{\log(Q/P)}{\log(2B+1)}.
\end{equation}
\end{theorem}

\begin{proof}
Write $c_a=-\log w_a\ge0$ and $\tau_d=-\log D_d$. Choose a pair
$(a_d,b_d)$ attaining $D_d$ for each positive $d\in\D$, and set
$(a_{-d},b_{-d})=(b_d,a_d)$. For $d=0$, use $(a_*,a_*)$, where
$w_{a_*}=1$. Then $\tau_0=0$ and $\tau_{-d}=\tau_d$. Define
\[
 p_d=\frac{D_d}{Q},
 \qquad
 r=\frac12\sum_{d\in\D}p_d\tau_d,
 \qquad
 V_m=\left\{x\in\A^m:\sum_{i=1}^m c_{x_i}\le rm\right\}.
\]
The probabilities $p_d$ are positive and satisfy $p_{-d}=p_d$.

\smallskip\noindent\emph{Difference count.}
For $d\ne0$, take $n_d=\lfloor mp_d\rfloor$ and put
$n_0=m-\sum_{d\ne0}n_d$. Hence $n_d=n_{-d}$ and $n_d/m\to p_d$ for
every $d$. Consider any ordering $z\in\D^m$ of this multiset of
differences. Replace each coordinate $d$ by the selected pair
$(a_d,b_d)$, obtaining vectors $x,y$ with $x-y=z$.
Their two costs are equal, since opposite differences occur equally often.
Each cost is
\[
 \sum_{\substack{d\in\D\\d>0}}n_d(c_{a_d}+c_{b_d})
 \le m\sum_{\substack{d\in\D\\d>0}}p_d\tau_d
 =mr.
\]
The zero coordinates contribute no cost. Thus every such $z$ belongs to
$V_m-V_m$. The number of distinct orderings is
$m!/\prod_{d\in\D}n_d!$. Stirling's formula gives
\begin{align}
 \log|V_m-V_m|
 &\ge mH(p)-o(m),\label{eq:type}\\
 H(p)=-\sum_{d\in\D}p_d\log p_d
 &=\log Q+\sum_{d\in\D}p_d\tau_d
 =\log Q+2r.\label{eq:entropy}
\end{align}

\smallskip\noindent\emph{Sum count.}
If $z=x+y$ for $x,y\in V_m$, the defining maximum in~\eqref{eq:sums}
implies
\[
 \prod_{i=1}^m S_{z_i}
 \ge\prod_{i=1}^m w_{x_i}w_{y_i}
 \ge e^{-2rm}.
\]
Summing these nonnegative products over all of $\Sset^m$ yields $P^m$.
Consequently
\begin{equation}\label{eq:sum-count}
 |V_m+V_m|\le e^{2rm}P^m.
\end{equation}
Combining~\eqref{eq:type}--\eqref{eq:sum-count},
\begin{equation}\label{eq:ratio}
 \log\frac{|V_m-V_m|}{|V_m+V_m|}
 \ge m\log(Q/P)-o(m).
\end{equation}

\smallskip\noindent\emph{Encoding as integers.}
Set $L=2B+1$ and
$\phi_m(x)=\sum_{i=1}^m x_iL^{i-1}$. This map is injective on both
$\{0,\ldots,2B\}^m$ and $\{-B,\ldots,B\}^m$.
Indeed, if two different strings first differ at their highest index $j$,
that term has absolute value at least $L^{j-1}$, whereas all lower terms
together have absolute value at most
$2B\sum_{i<j}L^{i-1}=L^{j-1}-1$.

Let $U_m=\phi_m(V_m)-\min\phi_m(V_m)$. It contains zero, its sum and
difference cardinalities equal those of $V_m$, and
\begin{equation}\label{eq:diameter}
 2\max U_m+1
 \le 2B\sum_{i=1}^m L^{i-1}+1
 =L^m.
\end{equation}
For all sufficiently large $m$, the logarithm in~\eqref{eq:ratio} is
positive, so $M_m=\max U_m>0$. Set $W_m=\{2u:u\in U_m\}$.
This dilation preserves both sum and difference cardinalities and ensures
the strict hypothesis of~\eqref{eq:transfer}, because
\[
 |W_m-W_m|=|U_m-U_m|
 \le2M_m+1<4M_m+1=2\max W_m+1.
\]
Moreover,~\eqref{eq:diameter} gives
$2\max W_m+1\le2L^m-1$. Applying~\eqref{eq:transfer} to $W_m$ yields
\[
 \theta\ge
 1+\frac{m\log(Q/P)-o(m)}{m\log L+\log2}.
\]
Taking $m\to\infty$ proves~\eqref{eq:bound} for the supremum $\theta$.
In particular, every exponent strictly below its right-hand side is
attainable in~\eqref{eq:problem}; we do not assert attainment of the limit.
\end{proof}

\section{Sparse geometric weights}
For an alphabet containing zero and geometric weights $w_a=q^a$, with
$0<q<1$, the two totals have a
particularly simple form:
\begin{equation}\label{eq:sparse}
 P(q)=\sum_{s\in\A+\A}q^s,
 \qquad
 Q(q)=\sum_{d\in\A-\A}q^{e(d)},
 \qquad
 e(d)=\min_{a-b=d}(a+b).
\end{equation}
Thus both totals are polynomials with nonnegative integer coefficients.
The formula follows directly from $w_aw_b=q^{a+b}$ and the strict decrease
of $q^j$ with $j$.

The finite integer family is also explicit. Define
\begin{equation}\label{eq:radius}
 \rho(q)=\frac{qQ'(q)}{2Q(q)},
 \qquad
 V_m=\left\{x\in\A^m:\sum_{i=1}^m x_i\le\lfloor m\rho(q)\rfloor\right\}.
\end{equation}
The radius $r$ in Theorem~\ref{thm:weighted} is
$(-\log q)\rho(q)$, so~\eqref{eq:radius} gives precisely its sets $V_m$.
Multiply their base-$(2B+1)$ encodings by two, as in that proof, to obtain
the sets used in the transfer inequality. Zero already belongs to this
integer family, and $\rho(q)$ is rational whenever $q$ is rational.

\subsection{A one-digit omission}

Take $B\ge3$ and $\A_B=\{0,2,3,\ldots,B\}$. Its sumset is
$\{0,2,3,\ldots,2B\}$ and its difference set is $\{-B,\ldots,B\}$.
For $|d|\ge2$, the pair $(|d|,0)$ or its reversal minimizes $a+b$, so
$e(d)=|d|$. The difference $1$ first occurs as $3-2$, giving
$e(1)=e(-1)=5$, and $e(0)=0$. Therefore
\begin{equation}\label{eq:one-hole}
 P_B(q)=\sum_{j=0}^{2B}q^j-q,
 \qquad
 Q_B(q)=1+2q^5+2\sum_{j=2}^{B}q^j.
\end{equation}
The omitted digit removes one sum, while every difference remains
represented. This changes the weighted ratio used by
Theorem~\ref{thm:weighted}.

\begin{corollary}\label{cor:explicit}
The exponent in~\eqref{eq:problem} satisfies $\theta>1.1813$.
\end{corollary}
\begin{proof}
Set $B=11$ and $q=377/500$ in~\eqref{eq:one-hole}. Define the positive
integers
\begin{align}
 N_S&=\sum_{\substack{0\le j\le22\\j\ne1}}
              377^j500^{22-j},\label{eq:NS}\\
 N_D&=500^{22}+2\cdot377^5 500^{17}
                   +2\sum_{j=2}^{11}377^j500^{22-j}.\label{eq:ND}
\end{align}
Then $Q_{11}(q)/P_{11}(q)=N_D/N_S$. Exact integer exponentiation verifies
\begin{equation}\label{eq:certificate}
 N_D^{10000}>23^{1813}N_S^{10000}.
\end{equation}
The accompanying verifier checks this inequality with integer arithmetic.
Taking logarithms
in~\eqref{eq:certificate} and applying Theorem~\ref{thm:weighted} gives
\[
 \theta\ge1+\frac{\log(N_D/N_S)}{\log23}
 >1+\frac{1813}{10000}=1.1813.
\]
\end{proof}

\subsection{Omitting the digits 1, 2 and 5}
Let $B\ge8$ and $\mathcal C_B=\{0,3,4,6,7,\ldots,B\}$.
Its sumset is $\{0,\ldots,2B\}\setminus\{1,2,5\}$ and its difference
set is $\{-B,\ldots,B\}$. The three missing positive digits still occur
as differences: $1=4-3$, $2=6-4$ and $5=8-3$. These pairs minimize the sum
of their endpoints, giving $e(\pm1)=7$, $e(\pm2)=10$ and $e(\pm5)=11$.
Every other nonzero difference uses a digit and zero. Thus
\begin{align}
 P(q)&=1+q^3+q^4+\sum_{j=6}^{2B}q^j,\label{eq:three-hole-sums}\\
 Q(q)&=1+2\left(q^3+q^4+\sum_{j=6}^{B}q^j+q^7+q^{10}+q^{11}\right).
 \label{eq:three-hole-differences}
\end{align}

\begin{corollary}\label{cor:headline}
The exponent in~\eqref{eq:problem} satisfies $\theta>1.1855$.
\end{corollary}
\begin{proof}
Take $B=16$, $q=16/19$ and $T_j=16^j19^{32-j}$. By
\eqref{eq:three-hole-sums}--\eqref{eq:three-hole-differences},
$Q(q)/P(q)=M_D/M_S$, where
\begin{align}
 M_S&=T_0+T_3+T_4+\sum_{j=6}^{32}T_j,\label{eq:MS}\\
 M_D&=T_0+2\left(T_3+T_4+\sum_{j=6}^{16}T_j+T_7+T_{10}+T_{11}\right).
 \label{eq:MD}
\end{align}
The exact integer check below verifies
\begin{equation}\label{eq:main-certificate}
 M_D^{2000}>33^{371}M_S^{2000}.
\end{equation}
Consequently Theorem~\ref{thm:weighted} gives
\[
 \theta\ge1+\frac{\log(M_D/M_S)}{\log33}
 >1+\frac{371}{2000}=1.1855.
\]
\end{proof}

\section{Exact reproduction and scope}
The certificate can be checked with the following complete Python program.
It uses integer arithmetic and the standard library only.
\begin{verbatim}
from math import gcd

p, q, B = 16, 19, 16
term = lambda j: p**j * q**(2*B-j)
S = term(0) + term(3) + term(4) + sum(term(j) for j in range(6, 33))
D = term(0) + 2*(term(3) + term(4) + sum(term(j) for j in range(6, 17))
                + term(7) + term(10) + term(11))
g = gcd(S, D)
if (D//g)**2000 <= 33**371 * (S//g)**2000:
    raise ArithmeticError("The certificate failed")
print("Verified: theta > 1.1855")
\end{verbatim}
The accompanying verifier also constructs the maxima in
\eqref{eq:sums}--\eqref{eq:differences} directly from all ordered digit
pairs and compares them with the displayed polynomial formulas. The integer inequality,
rather than a rounded logarithm or a numerical optimizer, certifies the
stated decimal threshold.

The theorem holds for every fixed alphabet and every set of positive
weights satisfying $Q>P$. The selected parameters need not maximize its
right-hand side. Nor does a finite search over alphabets determine the
optimal exponent in~\eqref{eq:problem}. This note records the construction,
its proof and an exact certificate. It does not assert that the bound or
method is new, or that the parameters attain a current record.

\begin{thebibliography}{9}
\bibitem{ghr}
K. Gyarmati, F. Hennecart and I. Z. Ruzsa,
\emph{Sums and differences of finite sets},
Functiones et Approximatio Commentarii Mathematici
\textbf{37}(1) (2007), 175--186.
\url{https://gyarmatikati.web.elte.hu/publ/sumdiffv.pdf}.

\bibitem{zheng}
F. Zheng,
\emph{Sums and differences of sets: a further improvement over AlphaEvolve},
arXiv:2506.01896 (2025).
\url{https://arxiv.org/abs/2506.01896}.

\bibitem{linli}
H. Lin and S. Li,
\emph{Settling the optimal exponent relating sumsets and difference sets},
arXiv:2607.27199v1 (2026).
\url{https://arxiv.org/abs/2607.27199v1}.

\bibitem{hill}
J. Beck, P. V. Ogren and A. Kobren,
\emph{Hill Sampling for Test-Time Scaling: A Simple and Better Alternative
to Repeated Sampling, Evolution, and Training},
arXiv:2609.25510v2 (2026).
\url{https://arxiv.org/abs/2609.25510v2}.
\end{thebibliography}
\end{document}
